\documentclass[a4paper,12pt,twoside]{memoir}

% Castellano
\usepackage[spanish,es-tabla]{babel}
\selectlanguage{spanish}
\usepackage[utf8]{inputenc}
\usepackage[T1]{fontenc}
\usepackage{lmodern} % Scalable font
\usepackage{microtype}
\usepackage{placeins}

\RequirePackage{booktabs}
\RequirePackage[table]{xcolor}
\RequirePackage{xtab}
\RequirePackage{multirow}

% Links
\usepackage[colorlinks]{hyperref}
\hypersetup{
	allcolors = {red}
}

% Ecuaciones
\usepackage{amsmath}

% Rutas de fichero / paquete
\newcommand{\ruta}[1]{{\sffamily #1}}

% Párrafos
\nonzeroparskip


% Imagenes
\usepackage{graphicx}
\newcommand{\imagen}[2]{
	\begin{figure}[!h]
		\centering
		\includegraphics[width=0.9\textwidth]{#1}
		\caption{#2}\label{fig:#1}
	\end{figure}
	\FloatBarrier
}

\newcommand{\imagenflotante}[2]{
	\begin{figure}%[!h]
		\centering
		\includegraphics[width=0.9\textwidth]{#1}
		\caption{#2}\label{fig:#1}
	\end{figure}
}



% El comando \figura nos permite insertar figuras comodamente, y utilizando
% siempre el mismo formato. Los parametros son:
% 1 -> Porcentaje del ancho de página que ocupará la figura (de 0 a 1)
% 2 --> Fichero de la imagen
% 3 --> Texto a pie de imagen
% 4 --> Etiqueta (label) para referencias
% 5 --> Opciones que queramos pasarle al \includegraphics
% 6 --> Opciones de posicionamiento a pasarle a \begin{figure}
\newcommand{\figuraConPosicion}[6]{%
  \setlength{\anchoFloat}{#1\textwidth}%
  \addtolength{\anchoFloat}{-4\fboxsep}%
  \setlength{\anchoFigura}{\anchoFloat}%
  \begin{figure}[#6]
    \begin{center}%
      \Ovalbox{%
        \begin{minipage}{\anchoFloat}%
          \begin{center}%
            \includegraphics[width=\anchoFigura,#5]{#2}%
            \caption{#3}%
            \label{#4}%
          \end{center}%
        \end{minipage}
      }%
    \end{center}%
  \end{figure}%
}

%
% Comando para incluir imágenes en formato apaisado (sin marco).
\newcommand{\figuraApaisadaSinMarco}[5]{%
  \begin{figure}%
    \begin{center}%
    \includegraphics[angle=90,height=#1\textheight,#5]{#2}%
    \caption{#3}%
    \label{#4}%
    \end{center}%
  \end{figure}%
}
% Para las tablas
\newcommand{\otoprule}{\midrule [\heavyrulewidth]}
%
% Nuevo comando para tablas pequeñas (menos de una página).
\newcommand{\tablaSmall}[5]{%
 \begin{table}
  \begin{center}
   \rowcolors {2}{gray!35}{}
   \begin{tabular}{#2}
    \toprule
    #4
    \otoprule
    #5
    \bottomrule
   \end{tabular}
   \caption{#1}
   \label{tabla:#3}
  \end{center}
 \end{table}
}

%
% Nuevo comando para tablas pequeñas (menos de una página).
\newcommand{\tablaSmallSinColores}[5]{%
 \begin{table}[H]
  \begin{center}
   \begin{tabular}{#2}
    \toprule
    #4
    \otoprule
    #5
    \bottomrule
   \end{tabular}
   \caption{#1}
   \label{tabla:#3}
  \end{center}
 \end{table}
}

\newcommand{\tablaApaisadaSmall}[5]{%
\begin{landscape}
  \begin{table}
   \begin{center}
    \rowcolors {2}{gray!35}{}
    \begin{tabular}{#2}
     \toprule
     #4
     \otoprule
     #5
     \bottomrule
    \end{tabular}
    \caption{#1}
    \label{tabla:#3}
   \end{center}
  \end{table}
\end{landscape}
}

%
% Nuevo comando para tablas grandes con cabecera y filas alternas coloreadas en gris.
\newcommand{\tabla}[6]{%
  \begin{center}
    \tablefirsthead{
      \toprule
      #5
      \otoprule
    }
    \tablehead{
      \multicolumn{#3}{l}{\small\sl continúa desde la página anterior}\\
      \toprule
      #5
      \otoprule
    }
    \tabletail{
      \hline
      \multicolumn{#3}{r}{\small\sl continúa en la página siguiente}\\
    }
    \tablelasttail{
      \hline
    }
    \bottomcaption{#1}
    \rowcolors {2}{gray!35}{}
    \begin{xtabular}{#2}
      #6
      \bottomrule
    \end{xtabular}
    \label{tabla:#4}
  \end{center}
}

%
% Nuevo comando para tablas grandes con cabecera.
\newcommand{\tablaSinColores}[6]{%
  \begin{center}
    \tablefirsthead{
      \toprule
      #5
      \otoprule
    }
    \tablehead{
      \multicolumn{#3}{l}{\small\sl continúa desde la página anterior}\\
      \toprule
      #5
      \otoprule
    }
    \tabletail{
      \hline
      \multicolumn{#3}{r}{\small\sl continúa en la página siguiente}\\
    }
    \tablelasttail{
      \hline
    }
    \bottomcaption{#1}
    \begin{xtabular}{#2}
      #6
      \bottomrule
    \end{xtabular}
    \label{tabla:#4}
  \end{center}
}

%
% Nuevo comando para tablas grandes sin cabecera.
\newcommand{\tablaSinCabecera}[5]{%
  \begin{center}
    \tablefirsthead{
      \toprule
    }
    \tablehead{
      \multicolumn{#3}{l}{\small\sl continúa desde la página anterior}\\
      \hline
    }
    \tabletail{
      \hline
      \multicolumn{#3}{r}{\small\sl continúa en la página siguiente}\\
    }
    \tablelasttail{
      \hline
    }
    \bottomcaption{#1}
  \begin{xtabular}{#2}
    #5
   \bottomrule
  \end{xtabular}
  \label{tabla:#4}
  \end{center}
}



\definecolor{cgoLight}{HTML}{EEEEEE}
\definecolor{cgoExtralight}{HTML}{FFFFFF}

%
% Nuevo comando para tablas grandes sin cabecera.
\newcommand{\tablaSinCabeceraConBandas}[5]{%
  \begin{center}
    \tablefirsthead{
      \toprule
    }
    \tablehead{
      \multicolumn{#3}{l}{\small\sl continúa desde la página anterior}\\
      \hline
    }
    \tabletail{
      \hline
      \multicolumn{#3}{r}{\small\sl continúa en la página siguiente}\\
    }
    \tablelasttail{
      \hline
    }
    \bottomcaption{#1}
    \rowcolors[]{1}{cgoExtralight}{cgoLight}

  \begin{xtabular}{#2}
    #5
   \bottomrule
  \end{xtabular}
  \label{tabla:#4}
  \end{center}
}


















\graphicspath{ {./img/} }

% Capítulos
\chapterstyle{bianchi}
\newcommand{\capitulo}[2]{
	\setcounter{chapter}{#1}
	\setcounter{section}{0}
	\chapter*{#2}
	\addcontentsline{toc}{chapter}{#2}
	\markboth{#2}{#2}
}

% Apéndices
\renewcommand{\appendixname}{Apéndice}
\renewcommand*\cftappendixname{\appendixname}

\newcommand{\apendice}[1]{
	%\renewcommand{\thechapter}{A}
	\chapter{#1}
}

\renewcommand*\cftappendixname{\appendixname\ }

% Formato de portada
\makeatletter
\usepackage{xcolor}
\newcommand{\tutor}[1]{\def\@tutor{#1}}
\newcommand{\course}[1]{\def\@course{#1}}
\definecolor{cpardoBox}{HTML}{E6E6FF}
\def\maketitle{
  \null
  \thispagestyle{empty}
  % Cabecera ----------------
\noindent\includegraphics[width=\textwidth]{cabecera}\vspace{1cm}%
  \vfill
  % Título proyecto y escudo informática ----------------
  \colorbox{cpardoBox}{%
    \begin{minipage}{.8\textwidth}
      \vspace{.5cm}\Large
      \begin{center}
      \textbf{TFG del Grado en Ingeniería Informática}\vspace{.6cm}\\
      \textbf{\LARGE\@title{}}
      \end{center}
      \vspace{.2cm}
    \end{minipage}

  }%
  \hfill\begin{minipage}{.20\textwidth}
    \includegraphics[width=\textwidth]{escudoInfor}
  \end{minipage}
  \vfill
  % Datos de alumno, curso y tutores ------------------
  \begin{center}%
  {%
    \noindent\LARGE
    Presentado por \@author{}\\ 
    en Universidad de Burgos --- \@date{}\\
    Tutor: \@tutor{}\\
  }%
  \end{center}%
  \null
  \cleardoublepage
  }
\makeatother

\newcommand{\nombre}{Jorge Zamora Marqués} %%% cambio de comando

% Datos de portada
\title{Aplicación para la gestión del sonido en un dispositivo móvil}
\author{\nombre}
\tutor{Carlos López Nodal}
\date{\today}

\begin{document}

\maketitle


\newpage\null\thispagestyle{empty}\newpage


%%%%%%%%%%%%%%%%%%%%%%%%%%%%%%%%%%%%%%%%%%%%%%%%%%%%%%%%%%%%%%%%%%%%%%%%%%%%%%%%%%%%%%%%
\thispagestyle{empty}


\noindent\includegraphics[width=\textwidth]{cabecera}\vspace{1cm}

\noindent D. Carlos López Nodal, profesor del departamento de Ingeniería Civil, área de Lenguajes y Sistemas Informáticos.

\noindent Expone:

\noindent Que el alumno D. \nombre, con DNI 72895344-H, ha realizado el Trabajo final de Grado en Ingeniería Informática titulado Aplicación para la gestión del sonido en un dispositivo móvil. 

\noindent Y que dicho trabajo ha sido realizado por el alumno bajo la dirección del que suscribe, en virtud de lo cual se autoriza su presentación y defensa.

\begin{center} %\large
En Burgos, {\large \today}
\end{center}

\vfill\vfill\vfill

% Author and supervisor
\begin{minipage}{0.45\textwidth}
\begin{flushleft} %\large
Vº. Bº. del Tutor:\\[2cm]
D. Carlos López Nodal
\end{flushleft}
\end{minipage}
\hfill
\begin{minipage}{0.45\textwidth}
\begin{flushleft} %\large
Vº. Bº. del co-tutor:\\[2cm]
D. nombre co-tutor
\end{flushleft}
\end{minipage}
\hfill

\vfill

% para casos con solo un tutor comentar lo anterior
% y descomentar lo siguiente
%Vº. Bº. del Tutor:\\[2cm]
%D. nombre tutor


\newpage\null\thispagestyle{empty}\newpage




\frontmatter

% Abstract en castellano
\renewcommand*\abstractname{Resumen}
\begin{abstract}
Se desarrollará una aplicación móvil que permitirá la configuración de unos perfiles de audio los cuales serán establecidos en nuestro dispositivo móvil por una serie de eventos configurados por el usuario, como eventos manuales, periódicos, wifi, GPS y de calendario.
\end{abstract}

\renewcommand*\abstractname{Descriptores}
\begin{abstract}
Servicio, volumen, perfil, GPS, calendario, wifi, observador, concurrencia, Android \ldots
\end{abstract}

\clearpage

% Abstract en inglés
\renewcommand*\abstractname{Abstract}
\begin{abstract}
A mobile application will be developed that will allow the configuration of some audio profiles which will be established in our mobile device by a series of events configured by the user, such as manual, periodic, wifi, GPS and calendar events.
\end{abstract}

\renewcommand*\abstractname{Keywords}
\begin{abstract}
Service, volume, profile, GPS, calendar, wifi, observer, concurrency, Android.
\end{abstract}

\clearpage

% Indices
\tableofcontents

\clearpage

\listoffigures

\clearpage

\listoftables
\clearpage

\mainmatter
\capitulo{1}{Introducción}


En los días en los que vivimos, mas de cinco mil millones de personas usan en su vida diaria un SmarthPhone, con lo que a más de una le ha pasado y le pasará, que su terminal suene en situaciones indebidas. Lo que quiero desarrollar en este Trabajo de Fin de Grado es una aplicación para SmartPhone, más concreta mente para los que usan el Sistema Operativo Android, para evitar en lo más posible, el que el terminal suene en momentos indebidos, y con ello ahorrarnos alguna que otra vergüenza.


El principal problema que encontramos, es que llevamos encima los terminales, ya sea en el trabajo, durante el estudio, durante las horas de dormir ect. Lo que quiero conseguir es posibilitar el configurar una serie de condiciones para que el terminal cambie solo de estado de volumen, los cuales añadiremos nosotros.


Para facilitar el uso a los usuarios, podremos utilizar los eventos de los calendarios del móvil, así como nuestra situación GPS, por las conexiones Wifi, y por eventos que añadamos manualmente dentro de la aplicación.


Los elementos con los que podremos interactuar serán el volumen general del móvil, el volumen de la música y de la alarma, así como configurar si queremos vibración o no.

\clearpage

\capitulo{2}{Objetivos del proyecto}

\paragraph{Objetivo principal}


El primero objetivo que se busca es posibilitar la automatización del control del volumen de nuestro terminal móvil con una sencilla aplicación, y que con el mínimo esfuerzo el usuario pueda configurar distintas configuraciones de audio, para luego activarlas con unos eventos específicos, como eventos de calendario o la localización GPS.

\section{Objetivos técnicos del proyecto}


\begin{itemize}
	\item Se estudiarán las diferentes librerías para el control de calendarios, GPS y wifi.
	\item Se estudiará el funcionamiento de las aplicaciones Android así como el ciclo de vida de estas.
	\item El estudio de diferentes formas de serializar datos en dispositivos móviles.
	\item El estudio de las diferentes versiones de Android, así como de las ventajas que nos ofrece cada una.
	\item La generación de documentación con Latex.
\end{itemize}

\section{Objetivos propios de la aplicación}


\begin{itemize}
	\item Poder guardar configuraciones de sonido propias, para el uso de estas dentro de la aplicación.
	\item Poder generar eventos propios por el usuario, ya sean de forma periódica o para momentos concretos.
	\item Poder obtener la red wifi a la que el usuario está conectado con su dispositivo.
	\item Poder leer los eventos de los calendarios de un usuario.
	\item Poder obtener la posición GPS de nuestro dispositivo Android.
	\item Poder configurar ciertas particualridades de los eventos, como distancia de aproximación en el GPS, lectura de la actual red wifi o poder elegir si queremos los eventos de calendario que su búsqueda sea exacta o que contenga una palabra clave.
	\item Poder obtener una rápida ayuda desde la propia aplicación para utilización de esta.
	\item Poder saber cual es el evento activo en cualquier momento.
	\item Universalizar los layouts de la aplicación.
	\item Internacionalización de la aplicación.
\end{itemize}
\capitulo{3}{Conceptos teóricos}

\section{Aplicación móvil}

Una aplicación móvil es es una aplicación informática con el objetivo de ser ejecutada dentro de dispositivos móviles como un smartphone o una tablet.

\section{Android}

A continuación hablaremos de como funciona el Sistema Operativo de Android, el cual está basado en el núcleo de Linux y desarrollado para su ejecución en dispositivos móviles como SmartPones, tablets, SmartWatchs...
La última versión de este Sistema Operativo es la \textit{9.0 "Pie"}.

\subsection{Arquitectura de Android}
La arquitectura de Andorid está estructurado en cinco niveles, siendo estos las aplicaciones, los marcos de trabajo de las aplicaciones, las bilbiotecas, el runtime de Android y el núcleo de Linux.
\imagen{AndroidArquitectura}{Arquitectura de Android}

Aplicaciones: las aplicaciones base incluyen un cliente de correo electrónico, programa de SMS, calendario, mapas, navegador, contactos y otros. Todas las aplicaciones están escritas en lenguaje de programación Java.

Marco de trabajo de aplicaciones: los desarrolladores tienen acceso completo a los mismos APIs del framework usados por las aplicaciones base. La arquitectura está diseñada para simplificar la reutilización de componentes; cualquier aplicación puede publicar sus capacidades y cualquier otra aplicación puede luego hacer uso de esas capacidades (sujeto a reglas de seguridad del framework). Este mismo mecanismo permite que los componentes sean reemplazados por el usuario.
Bibliotecas: Android incluye un conjunto de bibliotecas de C/C++ usadas por varios componentes del sistema. Estas características se exponen a los desarrolladores a través del marco de trabajo de aplicaciones de Android; algunas son: System C library (implementación biblioteca C estándar), bibliotecas de medios, bibliotecas de gráficos, 3D y SQLite, entre otras.

Runtime de Android: Android incluye un set de bibliotecas base que proporcionan la mayor parte de las funciones disponibles en las bibliotecas base del lenguaje Java. Cada aplicación Android corre su propio proceso, con su propia instancia de la máquina virtual Dalvik. Dalvik ha sido escrito de forma que un dispositivo puede correr múltiples máquinas virtuales de forma eficiente. Dalvik ejecutaba hasta la versión 5.0 archivos en el formato Dalvik Executable (.dex), el cual está optimizado para memoria mínima. La Máquina Virtual está basada en registros y corre clases compiladas por el compilador de Java que han sido transformadas al formato.dex por la herramienta incluida "dx". Desde la versión 5.0 utiliza el ART, que compila totalmente al momento de instalación de la aplicación.

Núcleo Linux: Android depende de Linux para los servicios base del sistema como seguridad, gestión de memoria, gestión de procesos, pila de red y modelo de controladores. El núcleo también actúa como una capa de abstracción entre el hardware y el resto de la pila de software.\cite{Android}

\section{Entorno de desarrollo}
El entorno de desarrollo integrado es una aplicación informática que permite el desarrollo de software de una manera más sencilla. Los entornos de desarrollo suelensonstar de un editor de exto, herramientas de construcción automáticas y un debuger.
Algunos entornos cuentan con IntelliSense para auto completar texto y un compilador, para poder compilar el código generado con el entorno de desarrollo.


\section{Java}
Java es un lenguaje de programación orientado a objetos, concurrente y de propósito general, y buscaba la independencia de la implementación de las aplicaciones, para conseguir poder ejecutar un mismo código sobre diferentes dispositivos con tecnologías diferentes.

Que sea un lenguaje de programación orientado a objetos quiere decir que cada tipo de dato que tengamos en la aplicación tendrá que llevar sus operaciones específicas para ese dato, es decir, cada objeto tendrá asociado un comportamiento y un estado.
Al separar estos comportamientos y estados unos de otros, permite de manera más sencilla el reutilizar código de diferentes aplicaciones.

Para permitir la independencia del código sobre las plataformas Java al compilarse genera un código llamado bytecode que son instrucciones que se le darán a la máquina que corre Java. Es decir, que cuando instalemos Java tendremos que seleccionar el que sea compatible con nuestro dispositivo, pero a la hora de ejecutar código, al este correr por la máquina virtual de Java lo que permite es que no necesite de especificaciones concretas de cada sistema el código de nuestra aplicación.

Otra característica principal de Java es el recolector de basura, que impedirá tener objetos que no se estén utilizando en este momento. Para ello el programador podrá especificar el ciclo de vida de los objetos, para controlar cuando se crean y cuando se borran. Es muy importante el recolector de basura ya que si no podríamos tener problemas con la memoria de nuestro dispositivo, ya que cada instancia de un objeto consume parte de esta.

Los entornos en los que funciona Java son sistemas móviles y embebidos, navegadores web, servidores y aplicaciones de escritorio.\cite{Java}

\section{Servicios Android - Services}
Un Service es un componente de una aplicación que puede realizar operaciones de larga ejecución en segundo plano y que no proporciona una interfaz de usuario. Otro componente de la aplicación puede iniciar un servicio y continuará ejecutándose en segundo plano aunque el usuario cambie a otra aplicación. Además, un componente puede enlazarse con un servicio para interactuar con él e incluso realizar una comunicación entre procesos (IPC). Por ejemplo, un servicio puede manejar transacciones de red, reproducir música, realizar I/O de archivos o interactuar con un proveedor de contenido, todo en segundo plano.

Un servicio puede adoptar esencialmente dos formas:

Servicio iniciado
Un servicio está "iniciado" cuando un componente de aplicación (como una actividad) lo inicia llamando a startService(). Una vez iniciado, un servicio puede ejecutarse en segundo plano de manera indefinida, incluso si se destruye el componente que lo inició. Por lo general, un servicio iniciado realiza una sola operación y no devuelve un resultado al emisor. Por ejemplo, puede descargar o cargar un archivo a través de la red. Cuando la operación está terminada, el servicio debe detenerse por sí mismo.
Servicio de enlace
Un servicio es de "de enlace" cuando un componente de la aplicación se vincula a él llamando a bindService(). Un servicio de enlace ofrece una interfaz cliente-servidor que permite que los componentes interactúen con el servicio, envíen solicitudes, obtengan resultados e incluso lo hagan en distintos procesos con la comunicación entre procesos (IPC). Un servicio de enlace se ejecuta solamente mientras otro componente de aplicación está enlazado con él. Se pueden enlazar varios componentes con el servicio a la vez, pero cuando todos ellos se desenlazan, el servicio se destruye.

\section{Listener}

Los listener son clases que se encargan de controlar los eventos que se generan dentro de nuestra aplicación.
Estos listeners podrán manejar los eventos para su posterior procesado, pudiendo capturar datos e información de dichos eventos.

\subsection{GPS}

GPS son las siglas de sistema de posicionamiento global, el cual nos permite establecer la posición de nuestro dispositivo móvil.
Se generará un listener para los eventos GPS lo que se nos permitirá saber de manera periódica donde se encuentra nuestro dispositivo móvil.

\section{Concurrencia}
La concurrencia es una propiedad que permite al sistema el poder realizar trabajos de manera simultánea. En nuestro caso, la concurrencia se realizará a través de los diferentes hilos de ejecución de nuestra aplicación, como el hilo principal de la aplicación (Main Thread), nuestro observador GPS el cual mirará en todo momento si se hacen cambios en el GPS, y el hilo del observador genérico, el cual es un hilo que se lanza cada x tiempo.

El principal problema que se trató con el tema de la concurrencia fueron el bloqueo del hilo principal por medio del observador genérico, ya que al utilizar un Thread nohabía control sobre el bloqueo del hilo principal. Como alternativas se utilizaron los AsincThread, hilis asíncronos que nos permitían la ejecución de tareas en segundo plano, aunque al final se decidió la utilización del TimerTask, ya que está preparado para la repetición de acciones programadas con un lapso x de tiempo.
\capitulo{4}{Técnicas y herramientas}

\section{Herramientas utilizadas}
Las herramientas utilizadas durante el proceso del desarrollo del trabajo han sido, el IDE Android Studio para la generación de código de Java, utilizando las librerías de Android, y a su vez el desarrollo de las interfaces gráficas de nuestra aplicación. Para el control del código generado también utilizamos Códacy, un herramienta que nos califica el código de nuestra aplicación.
Para gestionar las versiones de nuestra aplicación hemos utilizado Git-Hub y para complementarlo y gestionar la planificación del proyecto hemos utilizado Zen-Hub.
Para realizar la documentación hemos utilizado \TeX{}Maker y PowerPoint.

\subsection{Android Studio}

Es un entorno de desarrollo para aplicaciones Android que fue desarrollado por Google en 2014. Utiliza software de IntelliJ IDEA y es gratuito gracias a la licencia de Apache 2.0.
Las plataformas que soporta son Windows, macOS y Linux.

\subsection{\TeX{}Maker}

Es un editor de texto gratuito desarrollado por Pascal Brachet que permite el desarrollo de documentos con La\TeX{}.

\subsection{PowerPoint}

Aplicación de Microsoft que nos permite realizar de manera sencilla presentaciones.

\subsection{Codacy}

Es una herramienta en la nube que nos realiza un estudio sobre la calidad de código que está en un repositorio Git, como GitHub. Cada vez que realicemos un commit el aplicativo se encargará de dar un diagnóstico de la calidad de nuestro código así como de indicarnos donde están los problemas de este.

\subsection{GitHub + ZenHub}
GitHub es una plataforma de desarrollo de aplicaciones. Como su propio nombre indica GitHub utiliza Git, que es un software de control de versiones.
Complementado a GitHub encontramos a ZenHub, que es una extensión para Google Chroome que nos permite el gestionar proyectos con metodología Ágil.

\section{Técnicas utilizadas}

Las técnicas utilizadas a la hora de desarrollar el proyecto han sido dos, a la hora de planificar el desarrollo del proyecto hemos utilizado la metodología Ágil, y a la hora de realizar la programación hemos utilizado la metodología de programación orientada a objetos.

\subsection{Metodología Ágil}

La metodología ágil es una metodología de desarrollo software que se basa en el desarrollar los proyectos de una forma iterativa e incremental, donde los requisitos y sus soluciones van cambiando según las necesidades de ese memento.

Llamaremos iteración a la toma de decisiones y al trabajo que se realiza durante un periodo estipulado de tiempo. Dentro de cada iteración veremos una fase de planificación, de análisis de requisitos, diseño, codificación, pruebas y documentación.

En la fase de planificación lo que se planificará el trabajo que se tendrá que desarrollar en esa iteración. En la fase del análisis de requisitos se verá desarollarán las funcionalidades anteriormente planificadas. En la fase de diseño se verá gestionará la forma de abordar dichos requisitos. Las partes de codificación, documentación y pruebas son más parte del equipo de trabajo del proyecto, y serán los encargados del desarrollo dentro del software de las especificaciones anteriores, teniendo que realizar también la documentación del código generado y de las pruebas asociadas a cada parte del código.\cite{agil}

\subsection{Orientación a objetos.}
La programación orientada a objetos es una metodología de programación que se basa en el tratamiento de objetos que son entidades que tienen un estado, y los cuales pueden realizar unos determinados comportamientos.
Los objetos se descompondrán en atributos y métodos, siendo los atributos los encargados de especificar los estados de los objetos, y los métodos que son los encargados de especificar los comportamientos que pueden realizar dichos objetos.

Las principales características de la programación orientada a objetos son la capacidad para abstraer las características esenciales de los objetos, que todos los elementos pueden considerarse pertenecientes a una misma unidad, permite la herencia entre objetos, es decir que un objeto herede métodos y atributos de otra clase padre, permite la recolección de basura, las propiedades de los objetos están protegidas a modificaciones externas, y se permite la descomposición en partes más pequeñas el código generado.\cite{orientacionobjetos}

\tablaSmall{Herramientas y tecnologías utilizadas en cada parte del proyecto}{l c c c c}{herramientasportipodeuso}
{ \multicolumn{1}{l}{Herramientas} & Aplicación & BD & Memoria \\}{ 
Java & X & X &\\
Librerías Android Java & X & &\\
Android Studio & X & X &\\
Realm & & X &\\
Git + Zen-Hub & X & X & X\\
Mik\TeX{} & & & X\\
PowerPoint{} & & & X\\
\TeX{}Maker & & & X\\
Orientación a objetos & X & &\\
Metodología Ágil del desarrollo & X & X & X\\
} 
\capitulo{5}{Aspectos relevantes del desarrollo del proyecto}

\section{Inicio del proyecto}

Lo que me motivó a la hora de desarrollar este proyecto fue la necesidad de permitirme de una manera rápida y sencilla poder configurar mi terminal móvil para evitar situaciones embarazosas o no deseadas en las que el móvil suena cuando no debia, o por el contrario no sonaba cuando estabas disponible.

Para ello decidí realizar una aplicación que permitiera a un usuario no experto de los dispositivos móviles el poder mantener gestionado en todo momento y de manera sencilla el sonido de su dispositivo.

\section{Metodologías}

Dentro de lo posible se intentó llevar a cabo el desarrollo del proyecto siguiendo una metodología ágil, siendo esta Scrum, intentando ir sacando de manera ágil el contenido del proyecto.
Aunque no ha sido conseguido, como forma educativa se ha intentado seguir lo máximo posible esta idea.

Para ello se generaron una serie de hitos a ir cumpliendo divididos en diferentes sprints, que se buscaba que fuera el contenido de trabajo que se podía generar en una semana. Se debería de haber ido quedando con el tutor para la revisión de estos sprints pero por diferentes motivos se fueron fui haciendo revisiones autocríticas para la revisión de estos a la vez que se quedaba con el tutor y se mantenia conversación por correo electrónico.
Cara la planificación de estos sprints se utilizó ZenHub, una herramienta que representa un tablero canvas.

\subsection{Persistencia de los datos}
Otra decisión importante que se tuvo que tomar una vez ya arrancado el proyecto y avanzado, fue el modo de guardar los datos que introducía el usuario, es decir, cada vez que abriéramos la aplicación tendría que poder obtener los datos que habían sido guardados previamente, ya fueran eventos o configuraciones de sonido.

Las opciones que contemplamos fuero SQLlite, una base de datos con un tamaño reducido, que se ejecuta con la propia aplicación y Realm, que es un sistema de gestión de bases de datos diseñado inicialmente para Android y IOS de código abierto.
Al final se decidió la utilización de la base de datos Realm ya que su utilización era bastante sencilla.

La única complicación, es que hay que generar el id por el que va cada clase, para ello generaremos una clase que se ejecute al principio de cada ejecución de la aplicación, y que se encargara de obtener el ID máximo de cada clase, para evitar problemas de conflictos de PrimaryKey.

\subsection{LocationManager}
Location Manager es una clase propia de Android que nos permite utilizar los servicios de localización del SO. Con esta clase lo que podremos hacer será obtener la localización por varios métodos, en este caso la localización será la que obtengamos del GPS del móvil. 

\subsection{Easy Content Providers}

Es un proyecto de Android que permite obtener de manera sencilla diferentes datos del dispositivo, como por ejemplo los calendarios, contactos, log de llamadas\ldots

Con esta herramienta hemos conseguido iterar sobre los diferentes calendarios que se encuentran en nuestro dispositivo android, para luego obtener los eventos, y saber si hay alguno activo en ese momento o no.

\subsection{WifiManager}
WifiManager es la clase Java de Android, que nos permite obtener la información del wifi de nuestro terminal. Se estuvo pensando si utilizar toda la lista de redes wifi disponibles, la red con la intensidad mayor o la red wifi a la que estaba conectada el dispositivo, y al final me decanté por la red wifi a la que está conectado el dispositivo, ya que puede dar la casualidad que haya alguna otra red disponible con el mismo nombre que alguna de las que estuvieran configuradas, y de esta manera nos aseguramos que la red wifi con la que trataremos será la que utilice el usuario del terminal.

\subsection{Navigation Drawer}
Para facilitar la experiencia de usuario se pensó en la utilización de este tipo de menú, ya que es muy cómodo de utilizar y no ocupa espacio dentro de la interfaz de trabajo del usuario.
\capitulo{6}{Trabajos relacionados}

\section{No molestar (Android)}
"No molestar" es una configuración en los ajustes de los terminales Android que permite evitar la llegada de llamadas y notificaciones al dispositivo. Se podrán restringir números de teléfono y las diferentes notificaciones como eventos, llamadas y mensajes.\cite{nomolestar}

\section{Netroken Volume Control}
Es una aplicación que nos permite controlar los diferentes tipos de volúmenes dentro de nuestro dispositivo Android, por defecto ya hay unos creados y nosotros podremos crear más. Con la versión de pago podremos utilizar la localización de nuestro dispositivo para activar un perfil de sonido, a demás de programar cambios de perfil.\cite{netroken}

\section{Go Bees}
Go Bees es una aplicación móvil Android creada por el compañero David Miguel de nuestra universidad.
La pongo como trabajo relaccionado ya que es un buen ejemplo de como generar una buena aplicación Android con un seguimiento correcto del proyecto, documentación etc. La pongo como trabajo relacionado ya que lo utilizo como material de referencia para realizar este proyecto.\cite{gobees}

\section{Comparativa y opinión}
Comparando la app Netroken con la del proyecto, vemos que la idea es muy similar en ambas, aunque la aplicación de Netroken gane puntos ya que puede usarse como un gestor de volumen convencional, pierde también parte de su encanto al tener micro pagos para el desbloqueo de todas sus funcionalidades como la de la localización GPS para activar diferentes perfiles.
Pensándolo fríamente, las dos aplicaciones se basan en la misma premisa, pero siguen sin ser lo mismo, ya que este TFG se fundamenta en el control automático de los perfiles de sonido, sin darnos la oportunidad de poder cambiar de perfil de sonido manualmente. Esto a su vez le puede quitar valor con respecto a la aplicación Netroken.

En comparación con la aplicación de no molestar de Android, se podría decir que No Molestar es similar en parte a la configuración de los eventos periódicos, pero con el añadido de que nos permite, no solo configurar el volumen, si no la interacción con otros elementos como el filtrado de las llamadas.
\capitulo{7}{Conclusiones y Líneas de trabajo futuras}

\section{Conclusiones}
Comentándolo desde la parte del desarrollo de la aplicación, Android me ha parecido una tecnología bastante interesante, y a su vez me ha resultado conocida, ya que al trabajar con xml para el desarrollo de interfaces gráficas y la utilización de Java como lenguaje de programación, me ha resultado más sencillo de aprender.

Por otra parte los objetivos que tenía con este proyecto los he cumplido satisfactoriamente, ya que he conseguido crear una aplicación Android, que automatice el cambio de volumen del móvil con los eventos anteriormente citados.

Como opinión sobre la documentación de Android, creo que está bastante bien, aunque sin duda alguna, lo que me ha ayudado en gran medida a desarrollar el proyecto ha sido un curso de Android que he estado realizando como preparación.

Como pega al desarrollo de aplicaciones en Android, no solo necesitarás que en tu emulador funcione, si no que para cada versión y dispositivo. Con lo cual hay que estar siempre actualizando y mejorando tus aplicaciones, o si no se quedarán obsoletas.

\section{Nuevas líneas de trabajo}
A continuación hablaremos de posibles y propuestas a tener en cuenta para seguir con la evolución de este proyecto y de su coste en tiempo y del valor añadido que daría dicha acción.

\subsection{Servidor Web}
Para conseguir una mayor eficacia a la hora de gestionar las configuraciones de sonido de una persona, podríamos tener un servidor web, el cual se encargara de almacenar los datos, permitiéndonos conectarnos con una cuenta.
Coste: 80 - 100.
Valor añadido: 50 -100;

\subsection{Desarrollo para IOS}
Se podría desarrollar la aplicación para que se pueda integrar a los sistemas Apple.
Coste: 100 - 100.
Valor añadido: 100 - 100;

\subsection{Testeo de la aplicación en diferentes dispositivos y versiones}
Para asegurar un correcto funcionamiento de la aplicación, en un futuro se deberán de realizar más pruebas automáticas para luego estas ser probadas en servicios de virtualización de dispositivos, a si de esta manera podremos encontrar posibles errores de funcionamiento.
Coste 50 - 100.
Valor añadido 80 - 100.

\subsection{Publicación de la App}
Se podría publicar la App en el Play Store para que cualquier persona pudiera disfrutar de esta aplicación.
Coste: 10 - 100.
Valor añadido: 100 - 100.

\subsection{Creación de un nuevo evento de "caminando - corriendo"}
Se podría implementar, mediante la utilización del giroscopio del dispositivo, un evento que fuera lanzado cuando el usuario está caminando o corriendo con el dispositivo y habiendo trascurrido un tiempo x desde el inicio del desplazamiento.
Coste: 60 - 100.
Valor añadido: 40 - 100.

\subsection{Añadir control sobre la vibración}
Normalmente el control del volumen siempre va asociado al control de la vibración, con lo que se podría añadir el control de esta en nuestros perfiles de sonido.
Coste: 10 - 100.
Valor añadido: 60 - 100.


\bibliographystyle{plain}
\bibliography{bibliografia}

\end{document}
